% ============================================================
% CHƯƠNG 3: PHƯƠNG PHÁP ĐỀ XUẤT (Phần 2/2 — Mục 3.4–3.7)
% ============================================================
% Tiếp nối chapter3_part1.tex (Mục 3.1–3.3)

\section{Module 3: Pipeline xử lý dữ liệu và tập đặc trưng đầu vào}
\label{sec:data_pipeline}
\label{sec:feature_selection}

Trong phiên bản huấn luyện cuối cùng, khóa luận không xem lựa chọn đặc trưng như một module độc lập, mà tích hợp trực tiếp vào pipeline xử lý dữ liệu. Lý do là mục tiêu chính của bước này không phải xây dựng một thuật toán feature selection mới, mà là xác định một tập đầu vào gọn, dễ giải thích, phù hợp với bài toán dự đoán stress theo chuỗi thời gian và hạn chế nguy cơ rò rỉ dữ liệu.

Pipeline cuối cùng sử dụng \textbf{13 đặc trưng đầu vào}. Các đặc trưng này được chọn dựa trên ba tiêu chí:
\begin{itemize}
    \item \textbf{Tính liên quan với stress}: bao phủ các nhóm thời gian, hoạt động, cảm biến, sinh lý, hành vi và ngữ cảnh.
    \item \textbf{Tính khả dụng khi triển khai}: ưu tiên các trường có thể thu thập hoặc suy ra từ smartphone/wearable và nhật ký ngữ cảnh đơn giản.
    \item \textbf{An toàn cho dữ liệu chuỗi thời gian}: không sử dụng các đặc trưng rolling/aggregate có thể vô tình chứa thông tin tương lai.
\end{itemize}

\subsection{Tập đặc trưng đầu vào sử dụng cho huấn luyện}
\label{subsec:input_feature_set}

Bảng~\ref{tab:13_features_final} trình bày 13 đặc trưng được sử dụng làm đầu vào cho mô hình dự đoán stress.

\begin{table}[H]
\centering
\caption{Tập 13 đặc trưng đầu vào sử dụng cho huấn luyện}
\label{tab:13_features_final}
\small
\begin{tabular}{|l|c|p{8.5cm}|}
\hline
\textbf{Nhóm} & \textbf{Số lượng} & \textbf{Đặc trưng} \\
\hline
Thời gian (Temporal) & 2 & \texttt{Hour}, \texttt{Day\_of\_Week} \\
\hline
Hoạt động và cảm biến & 4 & \texttt{Activity}, \texttt{Accelerometer\_X}, \texttt{Accelerometer\_Y}, \texttt{Accelerometer\_Z} \\
\hline
Sinh lý (Physiological) & 1 & \texttt{Heart\_Rate} \\
\hline
Hành vi (Behavioral) & 3 & \texttt{Screen\_Usage\_Current}, \texttt{Phone\_Event\_Frequency}, \texttt{Mood\_Score} \\
\hline
Ngữ cảnh và phục hồi & 3 & \texttt{Location}, \texttt{Energy\_Level}, \texttt{Sleep\_Duration} \\
\hline
\textbf{Tổng} & \textbf{13} & -- \\
\hline
\end{tabular}
\end{table}

Ý nghĩa của từng nhóm đặc trưng:
\begin{itemize}
    \item \textbf{Thời gian}: \texttt{Hour} và \texttt{Day\_of\_Week} giúp mô hình học nhịp biến thiên stress theo giờ trong ngày và theo ngày trong tuần.
    \item \textbf{Hoạt động và cảm biến}: \texttt{Activity} cung cấp ngữ cảnh hoạt động, còn ba trục gia tốc kế giữ lại tín hiệu chuyển động thô tương ứng.
    \item \textbf{Sinh lý}: \texttt{Heart\_Rate} là chỉ số sinh lý trực tiếp, nhưng cần được diễn giải cùng ngữ cảnh hoạt động để tránh nhập nhằng.
    \item \textbf{Hành vi}: sử dụng điện thoại, tần suất sự kiện điện thoại và tâm trạng phản ánh các mẫu hành vi thường liên quan đến trạng thái căng thẳng.
    \item \textbf{Ngữ cảnh và phục hồi}: vị trí, năng lượng và thời lượng ngủ giúp mô hình nắm được môi trường hiện tại và khả năng phục hồi của người dùng.
\end{itemize}

Trong quá trình phát triển, khóa luận có khảo sát một tập đặc trưng ứng viên lớn hơn, bao gồm một số đặc trưng phái sinh và đặc trưng tổng hợp. Tuy nhiên, các đặc trưng dạng rolling hoặc aggregate có nguy cơ gây rò rỉ dữ liệu nếu không được xử lý rất chặt chẽ theo từng tập train/validation/test. Vì vậy, pipeline cuối cùng ưu tiên 13 đặc trưng trực tiếp, có khả năng diễn giải rõ và phù hợp hơn với mục tiêu xây dựng một baseline chuỗi thời gian đáng tin cậy.

Với cách trình bày này, tập đặc trưng đầu vào được xem là một phần của pipeline huấn luyện, không phải một đóng góp tách biệt. Trọng tâm của chương vì vậy được đặt vào quy trình xử lý dữ liệu chống rò rỉ và kiến trúc mô hình dự đoán stress.

\subsection{Vấn đề rò rỉ dữ liệu trong chuỗi thời gian}
\label{subsec:data_leakage_problem}

Data leakage (rò rỉ dữ liệu) xảy ra khi thông tin từ tập kiểm tra vô tình xuất hiện trong quá trình huấn luyện, dẫn đến hiệu suất đánh giá cao giả tạo \cite{kaufman2012leakage}. Trong dữ liệu chuỗi thời gian, data leakage đặc biệt nguy hiểm vì:

\begin{enumerate}
    \item \textbf{Không được xáo trộn (shuffle) dữ liệu}: Các mẫu liền kề có quan hệ thời gian chặt chẽ. Nếu shuffle, mẫu tương lai có thể xuất hiện trước mẫu hiện tại trong tập huấn luyện, khiến mô hình ``nhớ'' thay vì ``học'' xu hướng.
    \item \textbf{Chia dữ liệu phải theo thứ tự thời gian}: Train = phần đầu, Validation = phần giữa, Test = phần cuối — mô phỏng đúng tình huống triển khai thực tế.
    \item \textbf{Fit encoders/scalers chỉ trên train}: Nếu fit StandardScaler hoặc LabelEncoder trên toàn bộ dữ liệu, tham số thống kê (mean, std) sẽ chứa thông tin từ tập test.
\end{enumerate}

Khóa luận thiết kế pipeline xử lý dữ liệu tuân thủ nghiêm ngặt các nguyên tắc trên.

\subsection{Quy trình pipeline chi tiết}
\label{subsec:pipeline_detail}

Pipeline gồm 4 bước tuần tự, được minh họa trong Hình~\ref{fig:data_pipeline}:

\begin{figure}[H]
    \centering
    \includegraphics[width=0.9\linewidth]{chapter3/figs/data_pipeline.png}
    \caption{Pipeline xử lý dữ liệu chống rò rỉ}
    \label{fig:data_pipeline}
\end{figure}

\textbf{Bước 1 — Chia dữ liệu (Sequential Split):}

Dữ liệu 54.448 mẫu được chia theo thứ tự thời gian (không shuffle) với tỷ lệ 70/15/15:

\begin{equation}
\label{eq:data_split}
\begin{aligned}
\text{Train} &= \text{data}[0 : n_{\text{train}}] & \text{với } n_{\text{train}} &= \lfloor 0.70 \times N \rfloor = 38{,}113 \\
\text{Val}   &= \text{data}[n_{\text{train}} : n_{\text{val}}] & \text{với } n_{\text{val}} &= \lfloor 0.85 \times N \rfloor = 46{,}280 \\
\text{Test}  &= \text{data}[n_{\text{val}} : N] & & = 8{,}168 \text{ mẫu}
\end{aligned}
\end{equation}

Cách chia này đảm bảo tập test luôn nằm ở cuối dòng thời gian (tương ứng những ngày cuối cùng của 30 ngày mô phỏng), phản ánh đúng bài toán dự đoán tương lai.

\textbf{Bước 2 — Mã hóa biến phân loại (Label Encoding):}

Hai biến phân loại được mã hóa thành số nguyên bằng \texttt{LabelEncoder} của scikit-learn. Encoder được \textbf{fit chỉ trên tập train}, sau đó transform cho val và test:

\begin{table}[H]
\centering
\caption{Mã hóa biến phân loại}
\label{tab:label_encoding}
\small
\begin{tabular}{|l|c|l|}
\hline
\textbf{Biến} & \textbf{Số lớp} & \textbf{Giá trị} \\
\hline
Activity & 6 & Downstairs, Jogging, Sitting, Standing, Upstairs, Walking (mã hóa tự động theo LabelEncoder) \\
Location & 6 & commute, gym, home, outdoor, social, work (mã hóa tự động theo LabelEncoder) \\
\hline
\end{tabular}
\end{table}

\textbf{Bước 3 — Chuẩn hóa (Standard Normalization):}

Toàn bộ 13 đặc trưng (sau mã hóa) được chuẩn hóa bằng \texttt{StandardScaler}:

\begin{equation}
\label{eq:standardscaler}
x'_i = \frac{x_i - \mu_{\text{train}}}{\sigma_{\text{train}}}
\end{equation}

\noindent trong đó $\mu_{\text{train}}$ và $\sigma_{\text{train}}$ là giá trị trung bình và độ lệch chuẩn tính \textit{chỉ trên tập train}. Các tập val và test dùng cùng $\mu_{\text{train}}$, $\sigma_{\text{train}}$ — không fit lại.

\textbf{Bước 4 — Tạo chuỗi (Sequence Creation):}

Dữ liệu được biến đổi thành chuỗi bằng cửa sổ trượt (sliding window) với độ dài \texttt{seq\_length} = 60:

\begin{equation}
\label{eq:sequences}
\mathbf{X}_i = [\mathbf{x}_{i}, \mathbf{x}_{i+1}, \ldots, \mathbf{x}_{i+59}] \in \mathbb{R}^{60 \times 13}, \quad y_i = \text{Stress\_Level}_{i+60}
\end{equation}

\noindent Mỗi chuỗi gồm 60 mẫu liên tiếp (tương đương 30 phút ở tần suất 2 mẫu/phút), dùng để dự đoán giá trị stress tại thời điểm tiếp theo. Cửa sổ trượt di chuyển từng bước (stride = 1), tạo số lượng sequence lớn nhất có thể.

\begin{table}[H]
\centering
\caption{Kích thước dữ liệu qua từng bước pipeline}
\label{tab:pipeline_sizes}
\small
\begin{tabular}{|l|c|c|c|}
\hline
\textbf{Bước} & \textbf{Train} & \textbf{Val} & \textbf{Test} \\
\hline
1. Split (mẫu)     & 38.113 & 8.167 & 8.168 \\
2. Encode           & 38.113 $\times$ 13 & 8.167 $\times$ 13 & 8.168 $\times$ 13 \\
3. Normalize        & 38.113 $\times$ 13 & 8.167 $\times$ 13 & 8.168 $\times$ 13 \\
4. Sequences        & 38.053 $\times$ 60 $\times$ 13 & 8.107 $\times$ 60 $\times$ 13 & 8.108 $\times$ 60 $\times$ 13 \\
\hline
\end{tabular}
\end{table}

Số sequence = số mẫu $- $ \texttt{seq\_length} = $N - 60$. Ví dụ, tập train: $38.113 - 60 = 38.053$ sequences.

Điểm mấu chốt của pipeline: thứ tự \textbf{Split $\rightarrow$ Encode $\rightarrow$ Normalize $\rightarrow$ Sequences} đảm bảo không có rò rỉ dữ liệu ở bất kỳ bước nào. Nếu thực hiện encode/normalize trước split, hoặc tạo sequences trên toàn bộ dữ liệu rồi chia, thông tin thống kê từ tập test sẽ ảnh hưởng quá trình huấn luyện.

\subsection{Khả năng tái lập thực nghiệm}
\label{subsec:reproducibility}

Để tăng độ tin cậy khoa học và khả năng tái lập (reproducibility), khóa luận áp dụng các cơ chế kiểm soát sau trong toàn bộ pipeline huấn luyện:

\begin{itemize}
    \item \textbf{Cố định ngẫu nhiên:} sử dụng seed thống nhất (\texttt{RANDOM\_SEED = 42}) cho các thành phần ngẫu nhiên chính (NumPy, TensorFlow) trong các script huấn luyện baseline và tuning.
    \item \textbf{Chia dữ liệu tất định theo thời gian:} giữ thứ tự thời gian, không shuffle, với tỷ lệ 70/15/15 cố định để bảo đảm các lần chạy dùng cùng protocol đánh giá.
    \item \textbf{Tiền xử lý nhất quán:} tuân thủ nghiêm thứ tự Split $\rightarrow$ Encode $\rightarrow$ Normalize $\rightarrow$ Sequences; các encoder/scaler chỉ fit trên train và tái sử dụng cho val/test.
    \item \textbf{Lưu artifacts tiền xử lý và mô hình:} lưu \texttt{scaler}, \texttt{label encoders} và checkpoint mô hình tốt nhất, cho phép suy luận và phân tích lỗi trên cùng không gian đặc trưng đã huấn luyện.
\end{itemize}

Với cấu hình trên, các chỉ số MAE/RMSE/$R^2$ ổn định giữa các lần chạy lặp lại (dao động nhỏ do backend tính toán song song), đủ để hỗ trợ kết luận so sánh mô hình ở Chương~\ref{chap:experiments}.

% ============================================================

\section{Module 4: Kiến trúc mô hình Stacked Bidirectional LSTM}
\label{sec:model_architecture}

\subsection{Kiến trúc Baseline}
\label{subsec:baseline_architecture}

Mô hình baseline sử dụng kiến trúc Stacked Bidirectional LSTM gồm 2 lớp Bi-LSTM xếp chồng, kế thừa ưu điểm xử lý ngữ cảnh hai chiều đã chứng minh trong module HAR (Mục~\ref{sec:har_module}). Kiến trúc được thiết kế cho bài toán hồi quy (regression) dự đoán mức stress liên tục.

\begin{figure}[H]
    \centering
    \includegraphics[width=0.8\linewidth]{chapter3/figs/bilstm_baseline.png}
    \caption{Kiến trúc Stacked Bi-LSTM Baseline}
    \label{fig:bilstm_baseline}
\end{figure}

Chi tiết kiến trúc:

\begin{table}[H]
\centering
\caption{Kiến trúc chi tiết mô hình Baseline}
\label{tab:baseline_architecture}
\small
\begin{tabular}{|c|l|l|r|}
\hline
\textbf{Lớp} & \textbf{Tên lớp} & \textbf{Output Shape} & \textbf{Params} \\
\hline
1 & Input                              & (60, 13)         & 0 \\
2 & Bidirectional(LSTM(128, return\_seq=True)) & (60, 256)  & 145.408 \\
3 & Dropout(0.3)                       & (60, 256)        & 0 \\
4 & Bidirectional(LSTM(64))            & (128)            & 164.096 \\
5 & Dropout(0.3)                       & (128)            & 0 \\
6 & Dense(64, ReLU)                    & (64)             & 8.256 \\
7 & Dropout(0.3)                       & (64)             & 0 \\
8 & Dense(32, ReLU)                    & (32)             & 2.080 \\
9 & Dense(1, Linear)                   & (1)              & 33 \\
\hline
\multicolumn{3}{|r|}{\textbf{Tổng tham số}} & \textbf{320.129} \\
\hline
\end{tabular}
\end{table}

Giải thích thiết kế:
\begin{itemize}
    \item \textbf{Lớp Bi-LSTM đầu (128 units)}: Sử dụng \texttt{return\_sequences=True} để truyền toàn bộ hidden states cho lớp tiếp theo. Mỗi chiều (forward, backward) có 128 units, output tổng = 256 chiều. Lớp này nắm bắt các mẫu thời gian chi tiết (local patterns).
    \item \textbf{Lớp Bi-LSTM thứ hai (64 units)}: Chỉ trả về hidden state cuối cùng (mặc định \texttt{return\_sequences=False}), tổng hợp thông tin từ toàn bộ chuỗi thành vector 128 chiều. Lớp này nắm bắt xu hướng tổng thể (global trends).
    \item \textbf{Dropout(0.3)}: Regularization ngăn overfitting, ngẫu nhiên tắt 30\% neurons khi huấn luyện \cite{srivastava2014dropout}.
    \item \textbf{Dense layers (64 $\to$ 32 $\to$ 1)}: Giảm dần chiều (funnel architecture), với activation ReLU cho lớp ẩn. Lớp đầu ra dùng activation Linear phù hợp bài toán hồi quy.
\end{itemize}

Cấu hình huấn luyện baseline:
\begin{itemize}
    \item \textbf{Optimizer}: Adam \cite{kingma2014adam} với learning rate = 0.001.
    \item \textbf{Loss function}: Mean Squared Error (MSE).
    \item \textbf{Callbacks}:
    \begin{itemize}
        \item EarlyStopping(monitor=\texttt{val\_loss}, patience=15, restore\_best\_weights=True)
        \item ReduceLROnPlateau(monitor=\texttt{val\_loss}, factor=0.5, patience=5, min\_lr=$10^{-6}$)
        \item ModelCheckpoint — lưu mô hình tốt nhất theo val\_loss
    \end{itemize}
    \item \textbf{Batch size}: 64, \textbf{Max epochs}: 100.
\end{itemize}

\subsection{Tối ưu siêu tham số bằng Bayesian Optimization}
\label{subsec:bayesian_optimization}

Để cải thiện hiệu suất baseline, khóa luận sử dụng Bayesian Optimization \cite{snoek2012practical} thông qua thư viện Keras Tuner \cite{omalley2019kerastuner} để tìm kiếm tổ hợp siêu tham số tối ưu. So với Grid Search (thử tất cả tổ hợp) hay Random Search \cite{bergstra2012random} (thử ngẫu nhiên), Bayesian Optimization sử dụng mô hình xác suất (Gaussian Process) để dự đoán vùng tham số triển vọng, giúp tìm giải pháp tốt với ít trials hơn.

\textbf{Không gian tìm kiếm (Search Space):}

\begin{table}[H]
\centering
\caption{Không gian tìm kiếm siêu tham số}
\label{tab:search_space}
\small
\begin{tabular}{|l|l|l|}
\hline
\textbf{Siêu tham số} & \textbf{Kiểu} & \textbf{Phạm vi} \\
\hline
\texttt{lstm\_units\_1} (Bi-LSTM lớp 1) & Choice    & \{64, 128, 256\} \\
\texttt{lstm\_units\_2} (Bi-LSTM lớp 2) & Choice    & \{32, 64, 128\} \\
\texttt{dropout\_rate}                  & Float     & [0.1, 0.5], step 0.1 \\
\texttt{dense\_units}                   & Choice    & \{32, 64, 128\} \\
\texttt{learning\_rate}                 & Log Float & [$10^{-4}$, $10^{-2}$] \\
\hline
\end{tabular}
\end{table}

Tổng không gian: $3 \times 3 \times 5 \times 3 \times \text{continuous} = $ hàng trăm tổ hợp rời rạc, cộng thêm chiều liên tục của learning rate. Bayesian Optimization chỉ cần 20 trials (với 5 trials khám phá ngẫu nhiên ban đầu) để tìm giải pháp gần tối ưu.

\textbf{Cấu hình tuning:}
\begin{itemize}
    \item Số trials tối đa: 20 (5 random + 15 Bayesian)
    \item Epochs mỗi trial: 30 (với EarlyStopping patience=7)
    \item Objective: minimize \texttt{val\_mae}
    \item Batch size: 32
    \item Retraining: Mô hình tốt nhất được huấn luyện lại với 80 epochs và patience=15
\end{itemize}

\subsection{Kiến trúc Tuned (sau tối ưu)}
\label{subsec:tuned_architecture}

Bayesian Optimization tìm được tổ hợp siêu tham số tối ưu sau 20 trials. Bảng~\ref{tab:baseline_vs_tuned} so sánh cấu hình baseline và tuned.

\begin{table}[H]
\centering
\caption{So sánh cấu hình Baseline và Tuned}
\label{tab:baseline_vs_tuned}
\small
\begin{tabular}{|l|c|c|l|}
\hline
\textbf{Tham số} & \textbf{Baseline} & \textbf{Tuned} & \textbf{Nhận xét} \\
\hline
\texttt{lstm\_units\_1}  & 128 & \textbf{64}   & Giảm 50\% — mô hình nhỏ hơn \\
\texttt{lstm\_units\_2}  & 64  & \textbf{64}   & Giữ nguyên \\
\texttt{dropout\_rate}   & 0.3 & \textbf{0.1}  & Giảm mạnh — ít regularization hơn \\
\texttt{dense\_units}    & 64  & \textbf{128}  & Tăng — mở rộng lớp Dense \\
\texttt{learning\_rate}  & 0.001 & \textbf{0.01} & Tăng 10$\times$ — hội tụ nhanh hơn \\
\hline
\textbf{Tổng params}     & 320.129 & \textbf{163.585} & Giảm 48.9\% \\
\hline
\end{tabular}
\end{table}

Kiến trúc mô hình tuned:

\begin{table}[H]
\centering
\caption{Kiến trúc chi tiết mô hình Tuned}
\label{tab:tuned_architecture}
\small
\begin{tabular}{|c|l|l|r|}
\hline
\textbf{Lớp} & \textbf{Tên lớp} & \textbf{Output Shape} & \textbf{Params} \\
\hline
1 & Input                              & (60, 13)        & 0 \\
2 & Bidirectional(LSTM(64, return\_seq=True)) & (60, 128) & 39.936 \\
3 & Dropout(0.1)                       & (60, 128)       & 0 \\
4 & Bidirectional(LSTM(64))            & (128)           & 98.816 \\
5 & Dropout(0.1)                       & (128)           & 0 \\
6 & Dense(128, ReLU)                   & (128)           & 16.512 \\
7 & Dropout(0.1)                       & (128)           & 0 \\
8 & Dense(64, ReLU)                    & (64)            & 8.256 \\
9 & Dense(1, Linear)                   & (1)             & 65 \\
\hline
\multicolumn{3}{|r|}{\textbf{Tổng tham số}} & \textbf{163.585} \\
\hline
\end{tabular}
\end{table}

\textbf{Phân tích thay đổi:}
\begin{itemize}
    \item \textbf{Dropout 0.3 $\to$ 0.1}: Baseline có dấu hiệu underfitting nhẹ (val\_loss vẫn giảm khi dừng). Giảm dropout cho phép mô hình sử dụng nhiều neurons hơn trong quá trình huấn luyện, tăng khả năng biểu diễn.
    \item \textbf{LSTM units 128 $\to$ 64}: Giảm kích thước lớp Bi-LSTM đầu tiên nhưng giữ lớp thứ hai tại 64, tạo kiến trúc ``song song'' (equal-width) thay vì ``phễu'' (funnel) — phù hợp hơn khi dữ liệu không quá phức tạp.
    \item \textbf{Dense 64 $\to$ 128}: Mở rộng lớp Dense bù đắp cho việc giảm LSTM units, tăng khả năng mapping từ biểu diễn sequence sang stress level.
    \item \textbf{Learning rate 0.001 $\to$ 0.01}: Tốc độ học cao hơn giúp mô hình thoát khỏi local minima nhanh hơn, kết hợp với ReduceLROnPlateau sẽ tự động giảm dần khi cần thiết.
\end{itemize}

Kết quả so sánh sẽ được trình bày chi tiết trong Chương~\ref{chap:experiments}.

% ============================================================

\section{Các mô hình so sánh}
\label{sec:comparison_models}

Để đánh giá hiệu quả của kiến trúc Stacked Bi-LSTM đề xuất, khóa luận so sánh với 3 mô hình khác, từ đơn giản đến phức tạp, tất cả sử dụng cùng pipeline xử lý dữ liệu (Mục~\ref{sec:data_pipeline}) để đảm bảo tính công bằng.

\subsection{MLP (Multi-Layer Perceptron)}
\label{subsec:mlp_model}

Mô hình MLP là baseline đơn giản nhất, không có khả năng xử lý thông tin tuần tự. Toàn bộ chuỗi đầu vào $(60 \times 13)$ được làm phẳng (flatten) thành một vector 780 chiều, sau đó đi qua các lớp Dense:

\begin{table}[H]
\centering
\caption{Kiến trúc MLP}
\label{tab:mlp_architecture}
\small
\begin{tabular}{|c|l|l|r|}
\hline
\textbf{Lớp} & \textbf{Tên lớp} & \textbf{Output} & \textbf{Params} \\
\hline
1 & Input          & (60, 13)  & 0 \\
2 & Flatten        & (780)     & 0 \\
3 & Dense(256, ReLU) & (256)   & 199.936 \\
4 & Dropout(0.3)   & (256)     & 0 \\
5 & Dense(128, ReLU) & (128)   & 32.896 \\
6 & Dropout(0.3)   & (128)     & 0 \\
7 & Dense(64, ReLU)  & (64)    & 8.256 \\
8 & Dropout(0.2)   & (64)      & 0 \\
9 & Dense(1, Linear) & (1)     & 65 \\
\hline
\multicolumn{3}{|r|}{\textbf{Tổng tham số}} & \textbf{241.153} \\
\hline
\end{tabular}
\end{table}

MLP xem toàn bộ chuỗi như một vector phẳng — mất hoàn toàn thông tin thứ tự thời gian. Mục đích so sánh: đánh giá mức độ đóng góp của thông tin tuần tự (temporal patterns) vào dự đoán stress.

\subsection{Simple LSTM (1 lớp, một chiều)}
\label{subsec:simple_lstm_model}

Simple LSTM sử dụng 1 lớp LSTM đơn (unidirectional), chỉ xử lý dữ liệu theo một chiều thời gian (quá khứ $\to$ hiện tại):

\begin{table}[H]
\centering
\caption{Kiến trúc Simple LSTM}
\label{tab:simple_lstm_architecture}
\small
\begin{tabular}{|c|l|l|r|}
\hline
\textbf{Lớp} & \textbf{Tên lớp} & \textbf{Output} & \textbf{Params} \\
\hline
1 & Input           & (60, 13)  & 0 \\
2 & LSTM(128)       & (128)     & 72.704 \\
3 & Dropout(0.3)    & (128)     & 0 \\
4 & Dense(64, ReLU) & (64)      & 8.256 \\
5 & Dropout(0.2)    & (64)      & 0 \\
6 & Dense(32, ReLU) & (32)      & 2.080 \\
7 & Dense(1, Linear) & (1)     & 33 \\
\hline
\multicolumn{3}{|r|}{\textbf{Tổng tham số}} & \textbf{83.073} \\
\hline
\end{tabular}
\end{table}

So với Stacked Bi-LSTM, Simple LSTM có \textit{ít tham số nhất} (83K) vì chỉ có 1 lớp LSTM, 1 chiều. Mục đích: đánh giá lợi ích khi thêm lớp thứ hai (stacking) và xử lý hai chiều (bidirectional).

\subsection{Stacked Bi-GRU}
\label{subsec:bigru_model}

GRU (Gated Recurrent Unit) \cite{cho2014learning} là biến thể đơn giản hơn LSTM, chỉ có 2 cổng (reset và update) thay vì 3 cổng (forget, input, output). Kiến trúc Stacked Bi-GRU được thiết kế tương tự Stacked Bi-LSTM baseline để so sánh trực tiếp:

\begin{table}[H]
\centering
\caption{Kiến trúc Stacked Bi-GRU}
\label{tab:bigru_architecture}
\small
\begin{tabular}{|c|l|l|r|}
\hline
\textbf{Lớp} & \textbf{Tên lớp} & \textbf{Output} & \textbf{Params} \\
\hline
1 & Input                                   & (60, 13)   & 0 \\
2 & Bidirectional(GRU(128, return\_seq=True)) & (60, 256)  & 109.824 \\
3 & Dropout(0.3)                            & (60, 256)  & 0 \\
4 & Bidirectional(GRU(64))                  & (128)      & 123.648 \\
5 & Dropout(0.3)                            & (128)      & 0 \\
6 & Dense(64, ReLU)                         & (64)       & 8.256 \\
7 & Dropout(0.3)                            & (64)       & 0 \\
8 & Dense(32, ReLU)                         & (32)       & 2.080 \\
9 & Dense(1, Linear)                        & (1)        & 33 \\
\hline
\multicolumn{3}{|r|}{\textbf{Tổng tham số}} & \textbf{243.841} \\
\hline
\end{tabular}
\end{table}

GRU có ít tham số hơn LSTM (243K vs 320K cùng cấu hình) do ít cổng hơn. Mục đích: đánh giá liệu cấu trúc cổng đơn giản hơn có đủ cho bài toán dự đoán stress hay không.

\subsection{Tóm tắt 5 kiến trúc so sánh}
\label{subsec:model_summary}

Bảng~\ref{tab:five_models} tóm tắt 5 mô hình được so sánh trong khóa luận, theo trình tự từ baseline đơn giản đến mô hình đề xuất và biến thể tối ưu.

\begin{table}[H]
\centering
\caption{Tóm tắt 5 kiến trúc mô hình}
\label{tab:five_models}
\small
\begin{tabular}{|c|l|c|c|l|}
\hline
\textbf{\#} & \textbf{Mô hình} & \textbf{Params} & \textbf{Recurrent} & \textbf{Đặc điểm chính} \\
\hline
1 & MLP (Dense)                 & 241.153 & Không     & Flatten sequence, mất temporal \\
2 & Simple LSTM                 & 83.073  & 1 lớp, 1 chiều & Baseline RNN đơn giản \\
3 & Stacked Bi-LSTM (Baseline)  & 320.129 & 2 lớp, 2 chiều & Kiến trúc đề xuất \\
4 & Stacked Bi-GRU              & 243.841 & 2 lớp, 2 chiều & Thay LSTM bằng GRU \\
5 & Stacked Bi-LSTM (Tuned)     & 163.585 & 2 lớp, 2 chiều & Tối ưu bằng Bayesian Opt. \\
\hline
\end{tabular}
\end{table}

Mọi mô hình đều:
\begin{itemize}
    \item Sử dụng cùng pipeline dữ liệu (Split $\to$ Encode $\to$ Normalize $\to$ Sequences);
    \item Cùng tập train/val/test (đảm bảo công bằng);
    \item Optimizer Adam, loss MSE;
    \item Callbacks: EarlyStopping (patience=15), ReduceLROnPlateau (factor=0.5, patience=5);
    \item Batch size = 32, max epochs = 80.
\end{itemize}

\noindent Riêng mô hình Tuned (mô hình 5) có learning rate và dropout khác do kết quả tối ưu hóa. Kết quả so sánh chi tiết sẽ được trình bày trong Mục~\ref{sec:model_comparison_results} của Chương~\ref{chap:experiments}.
